\section{Related Work}
\label{sec:related_work}

Our proposed Lotka-Volterra Competitive Serving (LVCS) bridges parameter-efficient fine-tuning, stateful temporal serving, and biomimetic complex systems. In this section, we situate LVCS in the context of prior literature.

\subsection{Parameter-Efficient Fine-Tuning and Adapter Ensembling}
The growing scale of pre-trained models has made full parameter fine-tuning for multiple downstream tasks increasingly unsustainable. Low-Rank Adaptation (LoRA) \cite{hu2021lora} has emerged as the dominant PEFT paradigm, injecting trainable rank-decomposition matrices into attention and feed-forward layers while keeping the backbone frozen. Building upon LoRA, several works explore multi-task serving by deploying multiple specialized adapters on a shared backbone.

To route queries to these specialized adapters, early methods such as LoRA-Hub \cite{lora-hub}, LORS \cite{lors}, and MoFTE \cite{mofte} proposed static or query-level routing. More recently, SABLE \cite{sable} introduced a sample-wise nearest-centroid routing mechanism in the activation space of early layers. While SABLE represents a major milestone in stateless routing, it does not maintain any state across network depth or across temporal sequences. This absence of state results in a severe \emph{temporal routing jitter paradox}, where high-frequency sequential query noise causes routing coefficients to oscillate wildly from query to query over time. Furthermore, because SABLE lacks any active depth-wise recurrence, its ensembling weights remain highly soft and close to uniform across all layers, leading to representation leakage and co-dominance under overlapping task distributions. Our work addresses this bottleneck by introducing stateful, non-linear population dynamics across depth.

\subsection{Stateful Serving and Test-Time Adaptation}
Recognizing that real-world serving involves sequential query streams, several works have proposed stateful routing to enforce temporal and layer-by-layer smoothness. This line of research is closely related to Test-Time Adaptation (TTA) \cite{wang2020tent, zhang2022memo} and online adaptation under distribution shift \cite{shao2023test, lopez2017gradient}. 

ChemMerge \cite{chemmerge} introduced the concept of continuous-time biochemical reaction kinetics for model ensembling, modeling the routing state as concentrations of chemical species. However, ChemMerge is computationally expensive, requires ODE solvers, and relies on ad-hoc clipping to maintain weights on the probability simplex. Momentum-Merge \cite{momentummerge} simplified this process by demonstrating that continuous biochemistry collapses under Euler discretization to a simple Exponential Moving Average (EMA). PAC-Kinetics \cite{packinetics} formalized stateful routing by learning a linear state recurrence optimized via a PAC-Bayesian generalization bound for mixing stochastic processes.

Crucially, all of these stateful frameworks (ChemMerge, Momentum-Merge, PAC-Kinetics) rely on a linear state-space assumption. While linear state recurrences are mathematically convenient, they are fundamentally open-loop or linear-recurrent. They cannot capture the coupled, non-linear saturation and competitive dynamics required to model sharp task boundaries or stable co-existence under overlapping distributions. Furthermore, they suffer from representational lag during rapid task transitions. LVCS resolves these limitations by introducing non-linear Lotka-Volterra dynamics with Adaptive Niche Plasticity.

\subsection{Biomimetic AI and Ecological Modeling}
Drawing inspiration from biological systems has a rich history in computer science, from genetic algorithms \cite{holland1992adaptation} to neural networks themselves. In machine learning, researchers have applied biochemical ODEs to representational dynamics (e.g., Neural ODEs \cite{chen2018neural}) and evolutionary strategies for hyperparameter optimization \cite{jaderberg2017population}.

The Lotka-Volterra equations, originally proposed by Alfred Lotka \cite{lotka1925physical} and Vito Volterra \cite{volterra1926fluctuations}, serve as the mathematical foundation of theoretical ecology, modeling the dynamics of biological species competing for common resources. While Lotka-Volterra models have been used to analyze neural network competition \cite{fukushima1980neocognitron} and coordinate attention mechanisms in computer vision \cite{carandini2012normalization}, they have never been applied to the problem of stateful expert serving in multi-task deep networks. By formulating a discrete-time Lotka-Volterra Ricker competition recurrence \cite{ricker1954stock, may1976simple}, we establish a formal connection between mathematical ecology and dynamic ensembling, unlocking a new family of non-linear stateful routers.
